% Rainfall-Runoff Modeling Template
% Topics: Unit hydrograph theory, SCS Curve Number method, Nash cascade, watershed response
% Style: Engineering hydrology report with computational modeling

\documentclass[a4paper, 11pt]{article}
\usepackage[utf8]{inputenc}
\usepackage[T1]{fontenc}
\usepackage{amsmath, amssymb}
\usepackage{graphicx}
\usepackage{siunitx}
\usepackage{booktabs}
\usepackage{subcaption}
\usepackage[makestderr]{pythontex}
\usepackage{geometry}
\geometry{margin=1in}

% Theorem environments
\newtheorem{definition}{Definition}[section]
\newtheorem{theorem}{Theorem}[section]
\newtheorem{example}{Example}[section]
\newtheorem{remark}{Remark}[section]

\title{Rainfall-Runoff Modeling: Unit Hydrograph Theory and SCS Curve Number Method}
\author{Hydrology Research Group}
\date{\today}

\begin{document}
\maketitle

\begin{abstract}
This engineering report presents a comprehensive analysis of rainfall-runoff transformation
using the unit hydrograph approach and the Soil Conservation Service (SCS) Curve Number method.
We develop synthetic unit hydrographs using the Nash cascade model, compute excess rainfall
from design storms using the SCS-CN method, and apply discrete convolution to generate direct
runoff hydrographs. The analysis demonstrates watershed response characteristics including
time of concentration, peak discharge estimation, and the effects of antecedent moisture
conditions and land use on runoff generation.
\end{abstract}

\section{Introduction}

Rainfall-runoff modeling is fundamental to hydrologic engineering, providing the basis
for flood forecasting, drainage design, and water resources planning. The transformation
of precipitation into streamflow involves complex processes including interception,
infiltration, surface detention, and channel routing. The unit hydrograph approach,
introduced by Sherman (1932), provides an elegant framework for this transformation
by treating the watershed as a linear time-invariant system.

\begin{definition}[Unit Hydrograph]
A unit hydrograph (UH) is the direct runoff hydrograph resulting from 1 unit (typically
1 inch or 1 cm) of excess rainfall generated uniformly over the drainage area at a constant
rate for an effective duration $D$. The direct runoff excludes baseflow and represents
only the surface and rapid subsurface response.
\end{definition}

The SCS Curve Number method, developed by the USDA Natural Resources Conservation Service,
provides a practical approach to estimating rainfall excess from total precipitation based
on watershed characteristics including soil type, land use, and antecedent moisture conditions.

\section{Theoretical Framework}

\subsection{Unit Hydrograph Theory}

The unit hydrograph method relies on three fundamental principles:

\begin{enumerate}
\item \textbf{Time invariance}: The UH for a given watershed is time-invariant
\item \textbf{Proportionality}: Direct runoff ordinates are proportional to excess rainfall depth
\item \textbf{Superposition}: Runoff hydrographs from successive storms can be superposed
\end{enumerate}

\begin{theorem}[Convolution Equation]
For a watershed with unit hydrograph $u(t)$ and excess rainfall intensity $i_e(t)$,
the direct runoff $Q(t)$ is given by the convolution integral:
\begin{equation}
Q(t) = \int_0^t i_e(\tau) \cdot u(t - \tau) \, d\tau
\end{equation}
In discrete form with time intervals $\Delta t$:
\begin{equation}
Q_n = \sum_{m=1}^{n} P_m \cdot U_{n-m+1}
\end{equation}
where $P_m$ is excess rainfall in period $m$ and $U_n$ is the unit hydrograph ordinate.
\end{theorem}

\subsection{SCS Curve Number Method}

The SCS-CN method estimates direct runoff depth from rainfall depth based on the curve number $CN$,
which ranges from 0 (no runoff potential) to 100 (complete runoff).

\begin{definition}[SCS Runoff Equation]
The direct runoff depth $Q$ (inches) from rainfall depth $P$ (inches) is:
\begin{equation}
Q = \frac{(P - I_a)^2}{P - I_a + S}
\end{equation}
where $I_a$ is initial abstraction (interception + infiltration before runoff begins)
and $S$ is potential maximum retention after runoff begins. The relationships are:
\begin{equation}
S = \frac{1000}{CN} - 10 \quad \text{(inches)}
\end{equation}
\begin{equation}
I_a = 0.2S \quad \text{(standard assumption)}
\end{equation}
\end{definition}

\begin{remark}[Curve Number Selection]
Curve numbers depend on:
\begin{itemize}
\item \textbf{Hydrologic Soil Group}: A (low runoff), B, C, D (high runoff)
\item \textbf{Land Use/Cover}: Urban areas have higher CN than forests
\item \textbf{Antecedent Moisture Condition}: AMC I (dry), II (normal), III (wet)
\end{itemize}
\end{remark}

\subsection{Nash Cascade Model}

The Nash cascade represents the watershed as a series of $N$ identical linear reservoirs,
each with storage coefficient $K$.

\begin{theorem}[Nash Instantaneous Unit Hydrograph]
The instantaneous unit hydrograph (IUH) for a Nash cascade is:
\begin{equation}
u(t) = \frac{1}{K \Gamma(N)} \left(\frac{t}{K}\right)^{N-1} e^{-t/K}
\end{equation}
where $\Gamma(N)$ is the gamma function. The parameters relate to watershed properties:
\begin{equation}
K = \frac{t_p}{N} \quad \text{and} \quad N = \left(\frac{t_p}{\sigma}\right)^2
\end{equation}
where $t_p$ is time to peak and $\sigma$ is standard deviation of the IUH.
\end{theorem}

\section{Computational Analysis}

\subsection{Watershed Characteristics and Design Storm}

We analyze a small urban watershed to demonstrate rainfall-runoff transformation.
The watershed has mixed land use with moderate infiltration characteristics, typical
of suburban development with 40\% impervious cover. We apply a 6-hour design storm
with total precipitation depth of 4.5 inches, distributed using the SCS Type II
temporal distribution commonly used in the eastern United States.

\begin{pycode}
import numpy as np
import matplotlib.pyplot as plt
from scipy.special import gamma
from scipy.interpolate import interp1d

np.random.seed(42)

# Watershed characteristics
drainage_area = 2.5  # square miles
time_of_concentration = 1.2  # hours
CN = 75  # Curve number for suburban residential (1/4 acre lots, 38% impervious)

# Calculate storage parameters
S_inches = 1000/CN - 10  # Maximum retention (inches)
Ia_inches = 0.2 * S_inches  # Initial abstraction (inches)

# Design storm parameters
total_rainfall = 4.5  # inches
storm_duration = 6.0  # hours
dt = 0.25  # time step (hours)

# SCS Type II storm distribution (cumulative fractions)
# Time fractions of storm duration
time_fractions = np.array([0, 0.1, 0.2, 0.3, 0.35, 0.4, 0.45, 0.5, 0.55, 0.6,
                           0.65, 0.7, 0.75, 0.8, 0.85, 0.9, 0.95, 1.0])
# Cumulative rainfall fractions (Type II distribution)
cum_rainfall_fractions = np.array([0, 0.035, 0.076, 0.125, 0.156, 0.194, 0.219,
                                   0.663, 0.735, 0.772, 0.799, 0.820, 0.844,
                                   0.871, 0.898, 0.926, 0.963, 1.0])

# Create time series
time_storm = time_fractions * storm_duration
cum_rainfall = cum_rainfall_fractions * total_rainfall

# Interpolate to finer time step
time_series = np.arange(0, storm_duration + dt, dt)
interp_func = interp1d(time_storm, cum_rainfall, kind='linear', fill_value='extrapolate')
cum_rainfall_series = interp_func(time_series)
cum_rainfall_series = np.maximum(cum_rainfall_series, 0)
cum_rainfall_series = np.minimum(cum_rainfall_series, total_rainfall)

# Incremental rainfall depth
rainfall_incremental = np.diff(cum_rainfall_series, prepend=0)

# Calculate excess rainfall using SCS-CN method
excess_rainfall = np.zeros_like(cum_rainfall_series)
for i in range(len(cum_rainfall_series)):
    P = cum_rainfall_series[i]
    if P > Ia_inches:
        Q = ((P - Ia_inches)**2) / (P - Ia_inches + S_inches)
        excess_rainfall[i] = Q

# Incremental excess rainfall
excess_incremental = np.diff(excess_rainfall, prepend=0)

# Calculate runoff coefficient
total_excess = excess_rainfall[-1]
runoff_coefficient = total_excess / total_rainfall if total_rainfall > 0 else 0

# Plot rainfall and excess rainfall
fig1, (ax1, ax2) = plt.subplots(2, 1, figsize=(12, 8))

# Rainfall hyetograph
ax1.bar(time_series, rainfall_incremental / dt, width=dt*0.9,
        color='steelblue', edgecolor='black', alpha=0.7, label='Total rainfall')
ax1.bar(time_series, excess_incremental / dt, width=dt*0.9,
        color='darkred', edgecolor='black', alpha=0.7, label='Excess rainfall')
ax1.set_xlabel('Time (hours)', fontsize=11)
ax1.set_ylabel('Rainfall intensity (in/hr)', fontsize=11)
ax1.set_title(f'SCS Type II Design Storm (Total P = {total_rainfall} in, CN = {CN})', fontsize=12)
ax1.legend(fontsize=10, loc='upper left')
ax1.grid(True, alpha=0.3)
ax1.set_xlim(0, storm_duration)

# Cumulative rainfall and runoff
ax2.plot(time_series, cum_rainfall_series, 'b-', linewidth=2.5, label='Cumulative rainfall')
ax2.plot(time_series, excess_rainfall, 'r-', linewidth=2.5, label='Cumulative excess rainfall')
ax2.axhline(y=Ia_inches, color='gray', linestyle='--', linewidth=1.5,
            label=f'Initial abstraction = {Ia_inches:.2f} in')
ax2.fill_between(time_series, 0, excess_rainfall, color='darkred', alpha=0.2)
ax2.set_xlabel('Time (hours)', fontsize=11)
ax2.set_ylabel('Cumulative depth (inches)', fontsize=11)
ax2.set_title(f'Cumulative Rainfall and Runoff (Runoff coefficient = {runoff_coefficient:.3f})', fontsize=12)
ax2.legend(fontsize=10, loc='upper left')
ax2.grid(True, alpha=0.3)
ax2.set_xlim(0, storm_duration)
ax2.set_ylim(0, total_rainfall * 1.1)

plt.tight_layout()
plt.savefig('rainfall_runoff_scs_cn.pdf', dpi=150, bbox_inches='tight')
plt.close()
\end{pycode}

\begin{figure}[htbp]
\centering
\includegraphics[width=0.95\textwidth]{rainfall_runoff_scs_cn.pdf}
\caption{SCS Curve Number method applied to Type II design storm showing (top) rainfall
hyetograph with incremental rainfall intensity and excess rainfall intensity, and (bottom)
cumulative rainfall and runoff depths. The curve number CN = \py{CN} represents suburban
residential land use with moderate runoff potential. Initial abstraction $I_a$ = \py{f"{Ia_inches:.2f}"}
inches must be satisfied before runoff begins, resulting in a runoff coefficient of
\py{f"{runoff_coefficient:.3f}"}, indicating that \py{f"{runoff_coefficient*100:.1f}"}\%
of total precipitation becomes direct runoff. The nonlinear relationship between rainfall
and runoff is evident in the cumulative plot, with runoff generation accelerating once
initial abstraction is exceeded.}
\label{fig:scs_cn}
\end{figure}

\subsection{Synthetic Unit Hydrograph Development}

Using the time of concentration and drainage area, we develop a synthetic unit hydrograph
based on the Nash cascade model. The Nash model parameters are calibrated to match typical
watershed response characteristics for small urban basins. We use $N = 3$ reservoirs to
represent the storage-routing behavior of the watershed, which provides a realistic shape
with moderate skewness. The storage coefficient $K$ is derived from the time of concentration
using empirical relationships for urban watersheds.

\begin{pycode}
# Nash cascade parameters
N_reservoirs = 3  # Number of linear reservoirs
lag_time = 0.6 * time_of_concentration  # Lag time approximation

# Storage coefficient K from time to peak relationship
# For Nash model: tp ≈ K(N-1) for N > 1
K_hours = lag_time / (N_reservoirs - 1) if N_reservoirs > 1 else lag_time

# Time array for unit hydrograph (extended beyond storm)
time_uh = np.arange(0, 12, dt)

# Nash IUH function
def nash_iuh(t, K, N):
    """Nash instantaneous unit hydrograph"""
    if t <= 0:
        return 0
    return (1 / (K * gamma(N))) * ((t / K)**(N - 1)) * np.exp(-t / K)

# Generate instantaneous unit hydrograph
iuh = np.array([nash_iuh(t, K_hours, N_reservoirs) for t in time_uh])

# Convert IUH to unit hydrograph for duration D = dt
# UH(t) ≈ IUH(t) * dt for small dt
unit_hydrograph = iuh * dt

# Normalize to ensure sum equals 1 inch over drainage area
uh_sum = np.sum(unit_hydrograph)
if uh_sum > 0:
    unit_hydrograph = unit_hydrograph / uh_sum

# Convert to flow units (cfs per inch of excess rainfall)
# Q (cfs) = Depth (in) * Area (sq mi) * 640 acres/sq mi * 43560 sq ft/acre / (duration * 3600 s/hr)
conversion_factor = drainage_area * 640 * 43560 / (dt * 3600)
unit_hydrograph_cfs = unit_hydrograph * conversion_factor

# Calculate unit hydrograph properties
peak_uh = np.max(unit_hydrograph_cfs)
time_to_peak_uh = time_uh[np.argmax(unit_hydrograph_cfs)]
base_time_idx = np.where(unit_hydrograph_cfs > 0.01 * peak_uh)[0]
base_time_uh = time_uh[base_time_idx[-1]] if len(base_time_idx) > 0 else time_uh[-1]

# Plot unit hydrograph
fig2, ax = plt.subplots(figsize=(12, 7))
ax.fill_between(time_uh, 0, unit_hydrograph_cfs, color='steelblue', alpha=0.3, label='UH area')
ax.plot(time_uh, unit_hydrograph_cfs, 'b-', linewidth=2.5, label=f'Unit hydrograph (D = {dt} hr)')
ax.axvline(x=time_to_peak_uh, color='red', linestyle='--', linewidth=1.5,
           label=f'Time to peak = {time_to_peak_uh:.2f} hr')
ax.axhline(y=peak_uh, color='red', linestyle=':', linewidth=1.5, alpha=0.7)
ax.scatter([time_to_peak_uh], [peak_uh], s=100, c='red', zorder=5, edgecolor='black')

# Add annotations
ax.annotate(f'Peak = {peak_uh:.0f} cfs',
            xy=(time_to_peak_uh, peak_uh), xytext=(time_to_peak_uh + 1.5, peak_uh + 50),
            arrowprops=dict(arrowstyle='->', color='red', lw=1.5),
            fontsize=11, color='red', weight='bold')

ax.set_xlabel('Time (hours)', fontsize=12)
ax.set_ylabel('Discharge (cfs per inch of excess rainfall)', fontsize=12)
ax.set_title(f'Synthetic Unit Hydrograph - Nash Cascade Model (N = {N_reservoirs}, K = {K_hours:.2f} hr)',
             fontsize=13)
ax.legend(fontsize=11, loc='upper right')
ax.grid(True, alpha=0.3)
ax.set_xlim(0, 10)
ax.set_ylim(0, peak_uh * 1.2)

plt.tight_layout()
plt.savefig('rainfall_runoff_unit_hydrograph.pdf', dpi=150, bbox_inches='tight')
plt.close()
\end{pycode}

\begin{figure}[htbp]
\centering
\includegraphics[width=0.95\textwidth]{rainfall_runoff_unit_hydrograph.pdf}
\caption{Synthetic unit hydrograph derived from Nash cascade model with N = \py{N_reservoirs}
linear reservoirs and storage coefficient K = \py{f"{K_hours:.2f}"} hours. The unit hydrograph
represents the direct runoff response to 1 inch of excess rainfall uniformly distributed over
the \py{f"{drainage_area:.1f}"} square mile watershed during a \py{dt}-hour duration. Peak
discharge occurs at \py{f"{time_to_peak_uh:.2f}"} hours with magnitude \py{f"{peak_uh:.0f}"}
cfs per inch of excess rainfall. The base time is approximately \py{f"{base_time_uh:.1f}"}
hours, indicating the duration of direct runoff. The shape exhibits characteristic rising limb,
sharp peak, and gradual recession typical of small urban watersheds with rapid concentration
and moderate storage effects.}
\label{fig:unit_hydrograph}
\end{figure}

\subsection{Direct Runoff Hydrograph by Convolution}

We now apply discrete convolution of the excess rainfall hyetograph with the unit hydrograph
to generate the complete direct runoff hydrograph. This demonstrates the superposition principle,
where each increment of excess rainfall generates its own hydrograph following the unit hydrograph
shape, and all individual responses are summed to produce the total watershed response. The
convolution process accounts for the temporal distribution of rainfall intensity and the storage-
routing characteristics of the watershed.

\begin{pycode}
# Extend time series to accommodate full hydrograph
time_total = np.arange(0, storm_duration + base_time_uh + dt, dt)

# Extend excess rainfall array with zeros
excess_extended = np.zeros(len(time_total))
excess_extended[:len(excess_incremental)] = excess_incremental

# Extend unit hydrograph to match length
uh_extended = np.zeros(len(time_total))
uh_length = min(len(unit_hydrograph_cfs), len(uh_extended))
uh_extended[:uh_length] = unit_hydrograph_cfs[:uh_length]

# Perform discrete convolution
# Q(n) = sum(P(m) * U(n-m+1)) for m = 1 to n
direct_runoff = np.zeros(len(time_total))
for n in range(len(time_total)):
    for m in range(n + 1):
        if m < len(excess_extended) and (n - m) < len(uh_extended):
            direct_runoff[n] += excess_extended[m] * uh_extended[n - m]

# Calculate peak discharge and time to peak
peak_discharge = np.max(direct_runoff)
time_to_peak = time_total[np.argmax(direct_runoff)]

# Calculate runoff volume
runoff_volume_cf = np.sum(direct_runoff) * dt * 3600  # cubic feet
runoff_volume_inches = runoff_volume_cf / (drainage_area * 640 * 43560)

# Add baseflow (simplified constant baseflow)
baseflow = 15.0  # cfs
total_streamflow = direct_runoff + baseflow

# Calculate centroid of excess rainfall and runoff
excess_centroid = np.sum(time_series[:len(excess_incremental)] * excess_incremental) / np.sum(excess_incremental) if np.sum(excess_incremental) > 0 else 0
runoff_centroid = np.sum(time_total * direct_runoff * dt) / np.sum(direct_runoff * dt) if np.sum(direct_runoff * dt) > 0 else 0
lag_time_computed = runoff_centroid - excess_centroid

# Plot direct runoff hydrograph
fig3, (ax1, ax2) = plt.subplots(2, 1, figsize=(12, 10))

# Rainfall and runoff
ax1_twin = ax1.twinx()
ax1_twin.bar(time_series, rainfall_incremental / dt, width=dt*0.9,
             color='steelblue', alpha=0.4, edgecolor='black', linewidth=0.5)
ax1_twin.bar(time_series, excess_incremental / dt, width=dt*0.9,
             color='darkred', alpha=0.5, edgecolor='black', linewidth=0.7)
ax1_twin.set_ylabel('Rainfall intensity (in/hr)', fontsize=11, color='steelblue')
ax1_twin.tick_params(axis='y', labelcolor='steelblue')
ax1_twin.set_ylim(0, np.max(rainfall_incremental / dt) * 2.5)
ax1_twin.invert_yaxis()

ax1.fill_between(time_total, 0, direct_runoff, color='navy', alpha=0.3, label='Direct runoff')
ax1.plot(time_total, direct_runoff, 'b-', linewidth=2.5, label='Direct runoff hydrograph')
ax1.axvline(x=time_to_peak, color='red', linestyle='--', linewidth=1.5,
            label=f'Time to peak = {time_to_peak:.2f} hr')
ax1.axhline(y=peak_discharge, color='red', linestyle=':', linewidth=1.5, alpha=0.7)
ax1.scatter([time_to_peak], [peak_discharge], s=120, c='red', zorder=5,
            edgecolor='black', linewidth=1.5)
ax1.annotate(f'$Q_p$ = {peak_discharge:.0f} cfs',
             xy=(time_to_peak, peak_discharge), xytext=(time_to_peak + 1.5, peak_discharge * 0.85),
             arrowprops=dict(arrowstyle='->', color='red', lw=1.5),
             fontsize=12, color='red', weight='bold')

ax1.set_xlabel('Time (hours)', fontsize=11)
ax1.set_ylabel('Discharge (cfs)', fontsize=11)
ax1.set_title('Direct Runoff Hydrograph from Convolution of Excess Rainfall with Unit Hydrograph',
              fontsize=13)
ax1.legend(fontsize=10, loc='upper right')
ax1.grid(True, alpha=0.3)
ax1.set_xlim(0, storm_duration + base_time_uh)
ax1.set_ylim(0, peak_discharge * 1.3)

# Total streamflow with baseflow
ax2.fill_between(time_total, baseflow, total_streamflow, color='navy', alpha=0.3, label='Direct runoff')
ax2.fill_between(time_total, 0, baseflow, color='lightblue', alpha=0.5, label='Baseflow')
ax2.plot(time_total, total_streamflow, 'b-', linewidth=2.5, label='Total streamflow')
ax2.plot(time_total, np.ones_like(time_total) * baseflow, 'c--', linewidth=1.5, label=f'Baseflow = {baseflow} cfs')
ax2.axvline(x=time_to_peak, color='red', linestyle='--', linewidth=1.5, alpha=0.7)
ax2.scatter([time_to_peak], [peak_discharge + baseflow], s=120, c='red', zorder=5,
            edgecolor='black', linewidth=1.5)

ax2.set_xlabel('Time (hours)', fontsize=11)
ax2.set_ylabel('Discharge (cfs)', fontsize=11)
ax2.set_title(f'Total Streamflow Hydrograph (Direct Runoff + Baseflow)', fontsize=13)
ax2.legend(fontsize=10, loc='upper right')
ax2.grid(True, alpha=0.3)
ax2.set_xlim(0, storm_duration + base_time_uh)
ax2.set_ylim(0, (peak_discharge + baseflow) * 1.2)

plt.tight_layout()
plt.savefig('rainfall_runoff_hydrograph.pdf', dpi=150, bbox_inches='tight')
plt.close()
\end{pycode}

\begin{figure}[htbp]
\centering
\includegraphics[width=0.95\textwidth]{rainfall_runoff_hydrograph.pdf}
\caption{Direct runoff hydrograph generated by discrete convolution of excess rainfall
with the unit hydrograph (top panel) and complete streamflow hydrograph including baseflow
(bottom panel). The peak discharge $Q_p$ = \py{f"{peak_discharge:.0f}"} cfs occurs at
$t_p$ = \py{f"{time_to_peak:.2f}"} hours, demonstrating the lag between peak rainfall
intensity and peak watershed response. The computed lag time is \py{f"{lag_time_computed:.2f}"}
hours, measured between the centroids of excess rainfall and direct runoff. Total runoff
volume is \py{f"{runoff_volume_inches:.2f}"} inches, which matches the excess rainfall
calculated by the SCS-CN method, validating the convolution procedure. The hydrograph shape
exhibits rapid rise following intense rainfall, sustained peak during the main storm period,
and exponential recession characteristic of the Nash cascade routing. Baseflow contribution
of \py{baseflow} cfs represents the pre-storm groundwater contribution to streamflow.}
\label{fig:runoff_hydrograph}
\end{figure}

\section{Sensitivity Analysis}

\subsection{Effect of Curve Number on Runoff}

The curve number exerts strong control on runoff generation, with small changes in CN
producing significant changes in total runoff volume and peak discharge. We analyze
watershed response across a range of curve numbers representing different land use
conditions, from forested areas (low CN) to highly urbanized basins (high CN). This
sensitivity is particularly important for assessing the hydrologic impacts of land
development and urban sprawl.

\begin{pycode}
# Analyze sensitivity to curve number
CN_values = [60, 65, 70, 75, 80, 85, 90]  # Range from forest to urban
colors_cn = plt.cm.viridis(np.linspace(0.1, 0.9, len(CN_values)))

fig4, ((ax1, ax2), (ax3, ax4)) = plt.subplots(2, 2, figsize=(14, 10))

# Store results for each CN
peak_discharges_cn = []
runoff_coeffs = []
times_to_peak_cn = []

for i, CN_test in enumerate(CN_values):
    # Calculate excess rainfall for this CN
    S_test = 1000/CN_test - 10
    Ia_test = 0.2 * S_test

    excess_test = np.zeros_like(cum_rainfall_series)
    for j in range(len(cum_rainfall_series)):
        P = cum_rainfall_series[j]
        if P > Ia_test:
            Q = ((P - Ia_test)**2) / (P - Ia_test + S_test)
            excess_test[j] = Q

    excess_incr_test = np.diff(excess_test, prepend=0)

    # Convolution
    excess_ext_test = np.zeros(len(time_total))
    excess_ext_test[:len(excess_incr_test)] = excess_incr_test

    runoff_test = np.zeros(len(time_total))
    for n in range(len(time_total)):
        for m in range(n + 1):
            if m < len(excess_ext_test) and (n - m) < len(uh_extended):
                runoff_test[n] += excess_ext_test[m] * uh_extended[n - m]

    peak_q = np.max(runoff_test)
    peak_discharges_cn.append(peak_q)
    runoff_coeffs.append(excess_test[-1] / total_rainfall if total_rainfall > 0 else 0)
    times_to_peak_cn.append(time_total[np.argmax(runoff_test)])

    # Plot hydrograph
    ax1.plot(time_total, runoff_test, color=colors_cn[i], linewidth=2,
             label=f'CN = {CN_test}', alpha=0.8)

    # Plot cumulative runoff
    ax2.plot(time_total, np.cumsum(runoff_test) * dt * 3600 / (drainage_area * 640 * 43560),
             color=colors_cn[i], linewidth=2, label=f'CN = {CN_test}', alpha=0.8)

# Peak discharge vs CN
ax3.plot(CN_values, peak_discharges_cn, 'o-', color='darkblue', linewidth=2.5,
         markersize=8, markerfacecolor='steelblue', markeredgecolor='black', markeredgewidth=1.5)
ax3.set_xlabel('Curve Number (CN)', fontsize=11)
ax3.set_ylabel('Peak discharge (cfs)', fontsize=11)
ax3.set_title('Peak Discharge vs. Curve Number', fontsize=12)
ax3.grid(True, alpha=0.3)

# Runoff coefficient vs CN
ax4.plot(CN_values, runoff_coeffs, 's-', color='darkred', linewidth=2.5,
         markersize=8, markerfacecolor='coral', markeredgecolor='black', markeredgewidth=1.5)
ax4.set_xlabel('Curve Number (CN)', fontsize=11)
ax4.set_ylabel('Runoff coefficient', fontsize=11)
ax4.set_title('Runoff Coefficient vs. Curve Number', fontsize=12)
ax4.grid(True, alpha=0.3)
ax4.set_ylim(0, 1)

ax1.set_xlabel('Time (hours)', fontsize=11)
ax1.set_ylabel('Discharge (cfs)', fontsize=11)
ax1.set_title('Effect of Curve Number on Hydrograph Shape', fontsize=12)
ax1.legend(fontsize=9, loc='upper right')
ax1.grid(True, alpha=0.3)
ax1.set_xlim(0, 15)

ax2.set_xlabel('Time (hours)', fontsize=11)
ax2.set_ylabel('Cumulative runoff depth (inches)', fontsize=11)
ax2.set_title('Cumulative Direct Runoff Depth', fontsize=12)
ax2.legend(fontsize=9, loc='lower right')
ax2.grid(True, alpha=0.3)
ax2.set_xlim(0, 15)

plt.tight_layout()
plt.savefig('rainfall_runoff_cn_sensitivity.pdf', dpi=150, bbox_inches='tight')
plt.close()

# Calculate CN sensitivity metrics
CN_mid_idx = len(CN_values) // 2
dQ_dCN = (peak_discharges_cn[-1] - peak_discharges_cn[0]) / (CN_values[-1] - CN_values[0])
relative_change = (peak_discharges_cn[-1] - peak_discharges_cn[0]) / peak_discharges_cn[0] * 100
\end{pycode}

\begin{figure}[htbp]
\centering
\includegraphics[width=0.98\textwidth]{rainfall_runoff_cn_sensitivity.pdf}
\caption{Sensitivity of watershed response to curve number variation from CN = 60 (forested,
permeable soils) to CN = 90 (heavily urbanized, impervious surfaces). Top left panel shows
direct runoff hydrographs demonstrating dramatic increases in peak discharge and runoff volume
with increasing CN. Top right panel displays cumulative runoff depth, illustrating how
urbanization increases total runoff. Bottom panels quantify peak discharge sensitivity
(slope = \py{f"{dQ_dCN:.1f}"} cfs per CN unit) and runoff coefficient progression from
\py{f"{runoff_coeffs[0]:.2f}"} to \py{f"{runoff_coeffs[-1]:.2f}"}. The \py{f"{relative_change:.0f}"}\%
increase in peak discharge from CN = 60 to CN = 90 demonstrates the profound hydrologic impact
of urbanization, with corresponding increases in flood risk and reduced groundwater recharge.}
\label{fig:cn_sensitivity}
\end{figure}

\subsection{Effect of Time of Concentration}

Time of concentration affects both the unit hydrograph shape and watershed response timing.
Shorter times of concentration, typical of urbanized or steep watersheds, produce higher,
sharper peaks with reduced lag time. Longer times of concentration, characteristic of flat
rural watersheds with longer flow paths, produce attenuated hydrographs with delayed peaks.
We examine this sensitivity by varying the Nash model storage coefficient while maintaining
constant excess rainfall.

\begin{pycode}
# Analyze sensitivity to time of concentration
tc_values = [0.6, 0.8, 1.0, 1.2, 1.4, 1.6, 1.8]  # hours
colors_tc = plt.cm.plasma(np.linspace(0.1, 0.9, len(tc_values)))

fig5, ((ax1, ax2), (ax3, ax4)) = plt.subplots(2, 2, figsize=(14, 10))

peak_discharges_tc = []
times_to_peak_tc = []

for i, tc_test in enumerate(tc_values):
    # Recalculate unit hydrograph with new tc
    lag_test = 0.6 * tc_test
    K_test = lag_test / (N_reservoirs - 1) if N_reservoirs > 1 else lag_test

    iuh_test = np.array([nash_iuh(t, K_test, N_reservoirs) for t in time_uh])
    uh_test = iuh_test * dt
    uh_sum_test = np.sum(uh_test)
    if uh_sum_test > 0:
        uh_test = uh_test / uh_sum_test
    uh_test_cfs = uh_test * conversion_factor

    # Extend for convolution
    uh_ext_test = np.zeros(len(time_total))
    uh_length_test = min(len(uh_test_cfs), len(uh_ext_test))
    uh_ext_test[:uh_length_test] = uh_test_cfs[:uh_length_test]

    # Convolution with original excess rainfall
    runoff_tc_test = np.zeros(len(time_total))
    for n in range(len(time_total)):
        for m in range(n + 1):
            if m < len(excess_extended) and (n - m) < len(uh_ext_test):
                runoff_tc_test[n] += excess_extended[m] * uh_ext_test[n - m]

    peak_q_tc = np.max(runoff_tc_test)
    peak_discharges_tc.append(peak_q_tc)
    times_to_peak_tc.append(time_total[np.argmax(runoff_tc_test)])

    # Plot unit hydrograph
    ax1.plot(time_uh, uh_test_cfs, color=colors_tc[i], linewidth=2,
             label=f'$t_c$ = {tc_test} hr', alpha=0.8)

    # Plot runoff hydrograph
    ax2.plot(time_total, runoff_tc_test, color=colors_tc[i], linewidth=2,
             label=f'$t_c$ = {tc_test} hr', alpha=0.8)

# Peak discharge vs tc
ax3.plot(tc_values, peak_discharges_tc, 'o-', color='darkblue', linewidth=2.5,
         markersize=8, markerfacecolor='steelblue', markeredgecolor='black', markeredgewidth=1.5)
ax3.set_xlabel('Time of concentration (hours)', fontsize=11)
ax3.set_ylabel('Peak discharge (cfs)', fontsize=11)
ax3.set_title('Peak Discharge vs. Time of Concentration', fontsize=12)
ax3.grid(True, alpha=0.3)

# Time to peak vs tc
ax4.plot(tc_values, times_to_peak_tc, 's-', color='darkred', linewidth=2.5,
         markersize=8, markerfacecolor='coral', markeredgecolor='black', markeredgewidth=1.5)
ax4.plot(tc_values, tc_values, 'k--', linewidth=1.5, alpha=0.5, label='1:1 line')
ax4.set_xlabel('Time of concentration (hours)', fontsize=11)
ax4.set_ylabel('Time to peak discharge (hours)', fontsize=11)
ax4.set_title('Time to Peak vs. Time of Concentration', fontsize=12)
ax4.legend(fontsize=10)
ax4.grid(True, alpha=0.3)

ax1.set_xlabel('Time (hours)', fontsize=11)
ax1.set_ylabel('Discharge (cfs per inch)', fontsize=11)
ax1.set_title('Unit Hydrographs for Different $t_c$', fontsize=12)
ax1.legend(fontsize=9, loc='upper right')
ax1.grid(True, alpha=0.3)
ax1.set_xlim(0, 8)

ax2.set_xlabel('Time (hours)', fontsize=11)
ax2.set_ylabel('Discharge (cfs)', fontsize=11)
ax2.set_title('Runoff Hydrographs for Different $t_c$', fontsize=12)
ax2.legend(fontsize=9, loc='upper right')
ax2.grid(True, alpha=0.3)
ax2.set_xlim(0, 15)

plt.tight_layout()
plt.savefig('rainfall_runoff_tc_sensitivity.pdf', dpi=150, bbox_inches='tight')
plt.close()

# Calculate inverse relationship strength
from scipy.stats import linregress
slope_tc, intercept_tc, r_tc, p_tc, se_tc = linregress(tc_values, peak_discharges_tc)
\end{pycode}

\begin{figure}[htbp]
\centering
\includegraphics[width=0.98\textwidth]{rainfall_runoff_tc_sensitivity.pdf}
\caption{Sensitivity of watershed response to time of concentration $t_c$ ranging from
0.6 hours (highly urbanized, steep watershed) to 1.8 hours (rural, mild slopes). Top left
shows unit hydrographs becoming more attenuated and delayed with increasing $t_c$. Top right
displays corresponding runoff hydrographs exhibiting the inverse relationship between $t_c$
and peak discharge. Bottom left quantifies this relationship with regression slope
\py{f"{slope_tc:.1f}"} cfs/hour (R² = \py{f"{r_tc**2:.3f}"}), demonstrating that shorter
concentration times produce \py{f"{abs(peak_discharges_tc[0] - peak_discharges_tc[-1]):.0f}"}
cfs higher peaks. Bottom right shows time to peak increases approximately linearly with $t_c$,
though storm timing also influences this relationship. These results have critical implications
for urban drainage design and flood mitigation.}
\label{fig:tc_sensitivity}
\end{figure}

\section{Results Summary}

\begin{pycode}
print(r"\begin{table}[htbp]")
print(r"\centering")
print(r"\caption{Summary of Rainfall-Runoff Analysis Results}")
print(r"\begin{tabular}{lcc}")
print(r"\toprule")
print(r"\textbf{Parameter} & \textbf{Value} & \textbf{Units} \\")
print(r"\midrule")
print(r"\multicolumn{3}{l}{\textit{Watershed Characteristics}} \\")
print(f"Drainage area & {drainage_area:.2f} & sq mi \\\\")
print(f"Time of concentration & {time_of_concentration:.2f} & hr \\\\")
print(f"Curve number (CN) & {CN} & --- \\\\")
print(f"Maximum retention (S) & {S_inches:.2f} & in \\\\")
print(f"Initial abstraction ($I_a$) & {Ia_inches:.2f} & in \\\\")
print(r"\midrule")
print(r"\multicolumn{3}{l}{\textit{Design Storm}} \\")
print(f"Total rainfall depth & {total_rainfall:.2f} & in \\\\")
print(f"Storm duration & {storm_duration:.1f} & hr \\\\")
print(f"Distribution type & SCS Type II & --- \\\\")
print(r"\midrule")
print(r"\multicolumn{3}{l}{\textit{Runoff Generation}} \\")
print(f"Total excess rainfall & {total_excess:.2f} & in \\\\")
print(f"Runoff coefficient & {runoff_coefficient:.3f} & --- \\\\")
print(f"Runoff volume & {runoff_volume_cf:.0f} & cu ft \\\\")
print(r"\midrule")
print(r"\multicolumn{3}{l}{\textit{Unit Hydrograph}} \\")
print(f"Nash N parameter & {N_reservoirs} & --- \\\\")
print(f"Storage coefficient (K) & {K_hours:.2f} & hr \\\\")
print(f"Peak UH discharge & {peak_uh:.0f} & cfs/in \\\\")
print(f"Time to peak (UH) & {time_to_peak_uh:.2f} & hr \\\\")
print(f"Base time & {base_time_uh:.1f} & hr \\\\")
print(r"\midrule")
print(r"\multicolumn{3}{l}{\textit{Direct Runoff Hydrograph}} \\")
print(f"Peak discharge & {peak_discharge:.0f} & cfs \\\\")
print(f"Time to peak & {time_to_peak:.2f} & hr \\\\")
print(f"Lag time & {lag_time_computed:.2f} & hr \\\\")
print(r"\bottomrule")
print(r"\end{tabular}")
print(r"\label{tab:results}")
print(r"\end{table}")
\end{pycode}

\section{Discussion}

\subsection{Hydrologic Insights}

The analysis demonstrates several key principles of rainfall-runoff transformation:

\begin{enumerate}
\item \textbf{Nonlinear losses}: The SCS-CN method captures the nonlinear relationship
between precipitation and runoff, with runoff coefficient increasing from near-zero for
small storms to asymptotically approaching $(1 - I_a/P)$ for large storms.

\item \textbf{Watershed routing}: The Nash cascade model provides a parsimonious representation
of watershed storage and routing effects with physically meaningful parameters (N and K)
related to watershed geomorphology.

\item \textbf{Peak attenuation}: The convolution process demonstrates how watershed storage
attenuates and delays the runoff response relative to rainfall input, with peak discharge
occurring \py{f"{lag_time_computed:.2f}"} hours after the rainfall centroid.

\item \textbf{Land use sensitivity}: The CN sensitivity analysis quantifies the dramatic
impact of urbanization on flood potential, with peak discharge increasing by
\py{f"{relative_change:.0f}"}\% when CN increases from 60 to 90.
\end{enumerate}

\subsection{Engineering Applications}

These methods provide the foundation for:

\begin{itemize}
\item \textbf{Flood frequency analysis}: Design hydrographs for culvert and bridge sizing
\item \textbf{Stormwater management}: Detention basin sizing and outlet design
\item \textbf{Urban drainage}: Storm sewer network design and analysis
\item \textbf{Regulatory compliance}: Post-development runoff not exceeding pre-development conditions
\item \textbf{Watershed planning}: Assessing hydrologic impacts of land use changes
\end{itemize}

\subsection{Method Limitations}

Important assumptions and limitations include:

\begin{remark}[Model Assumptions]
\begin{itemize}
\item \textbf{Spatial uniformity}: Rainfall and watershed properties assumed spatially uniform
\item \textbf{Temporal invariance}: Unit hydrograph assumed constant for all storm events
\item \textbf{Linearity}: Superposition assumes linear watershed response
\item \textbf{Empiricism}: SCS-CN method is empirical, calibrated to agricultural watersheds
\item \textbf{Initial conditions}: Antecedent moisture significantly affects CN selection
\end{itemize}
\end{remark}

For large or heterogeneous watersheds, distributed hydrologic models may be more appropriate.

\section{Conclusions}

This comprehensive analysis of rainfall-runoff transformation demonstrates the coupled
application of the SCS Curve Number method and unit hydrograph theory for watershed response
prediction. Key findings include:

\begin{enumerate}
\item For the \py{f"{drainage_area:.1f}"} square mile urban watershed with CN = \py{CN},
the \py{total_rainfall}-inch design storm produces peak discharge $Q_p$ = \py{f"{peak_discharge:.0f}"}
cfs occurring at $t_p$ = \py{f"{time_to_peak:.2f}"} hours.

\item The runoff coefficient of \py{f"{runoff_coefficient:.3f}"} indicates that
\py{f"{runoff_coefficient*100:.1f}"}\% of total precipitation becomes direct runoff, with
\py{f"{Ia_inches:.2f}"} inches lost to initial abstraction and \py{f"{S_inches:.2f}"}
inches to continuing infiltration and retention.

\item The Nash cascade unit hydrograph with N = \py{N_reservoirs} reservoirs and K =
\py{f"{K_hours:.2f}"} hours accurately represents the storage-routing characteristics
of small urban watersheds, producing realistic hydrograph shape and timing.

\item Sensitivity analysis reveals peak discharge varies from \py{f"{peak_discharges_cn[0]:.0f}"}
to \py{f"{peak_discharges_cn[-1]:.0f}"} cfs across the CN range 60-90, emphasizing the
critical importance of land use in flood generation and the need for effective stormwater
management in developing areas.

\item Time of concentration exerts strong control on hydrograph attenuation, with the
inverse relationship between $t_c$ and peak discharge (slope = \py{f"{slope_tc:.1f}"}
cfs/hour) demonstrating why urbanization (which reduces $t_c$) amplifies flood peaks.
\end{enumerate}

These methods provide essential tools for hydrologic engineering practice, enabling
quantitative assessment of watershed response to design storms for infrastructure planning,
flood mitigation, and environmental protection.

\section*{References}

\begin{enumerate}
\item Sherman, L.K. (1932). Streamflow from rainfall by the unit-graph method.
\textit{Engineering News-Record}, 108, 501-505.

\item Soil Conservation Service (1985). \textit{National Engineering Handbook, Section 4: Hydrology}.
U.S. Department of Agriculture, Washington, D.C.

\item Chow, V.T., Maidment, D.R., and Mays, L.W. (1988). \textit{Applied Hydrology}.
McGraw-Hill, New York.

\item Nash, J.E. (1957). The form of the instantaneous unit hydrograph.
\textit{International Association of Scientific Hydrology, Publication}, 45(3), 114-121.

\item Hawkins, R.H., Ward, T.J., Woodward, D.E., and Van Mullem, J.A. (2009).
\textit{Curve Number Hydrology: State of the Practice}. American Society of Civil Engineers.

\item Natural Resources Conservation Service (2010). \textit{National Engineering Handbook,
Part 630: Hydrology}. U.S. Department of Agriculture.

\item Dooge, J.C.I. (1973). Linear theory of hydrologic systems. \textit{Technical Bulletin
No. 1468}, Agricultural Research Service, U.S. Department of Agriculture.

\item McCuen, R.H. (2016). \textit{Hydrologic Analysis and Design}, 4th ed. Pearson, Upper Saddle River, NJ.

\item Singh, V.P. (1988). \textit{Hydrologic Systems: Rainfall-Runoff Modeling}, Volume I.
Prentice Hall, Englewood Cliffs, NJ.

\item Viessman, W. and Lewis, G.L. (2003). \textit{Introduction to Hydrology}, 5th ed.
Pearson Education, Upper Saddle River, NJ.

\item Ponce, V.M. (1989). \textit{Engineering Hydrology: Principles and Practices}.
Prentice Hall, Englewood Cliffs, NJ.

\item USDA-NRCS (2004). Estimation of direct runoff from storm rainfall. In \textit{National
Engineering Handbook Part 630}, Chapter 10. U.S. Department of Agriculture.

\item Mishra, S.K. and Singh, V.P. (2003). \textit{Soil Conservation Service Curve Number (SCS-CN)
Methodology}. Kluwer Academic Publishers, Dordrecht.

\item Singh, V.P. (1992). \textit{Elementary Hydrology}. Prentice Hall, Englewood Cliffs, NJ.

\item Hjelmfelt, A.T. (1991). Investigation of curve number procedure. \textit{Journal of
Hydraulic Engineering}, 117(6), 725-737.

\item Mockus, V. (1949). Estimation of total (and peak rates of) surface runoff for individual
storms. \textit{Exhibit A of Appendix B, Interim Survey Report, Grand (Neosho) River Watershed},
U.S. Department of Agriculture.

\item Snyder, F.F. (1938). Synthetic unit-graphs. \textit{Transactions of the American Geophysical
Union}, 19, 447-454.

\item U.S. Army Corps of Engineers (2000). \textit{Hydrologic Modeling System HEC-HMS: Technical
Reference Manual}. Hydrologic Engineering Center, Davis, CA.

\item Singh, V.P. and Frevert, D.K. (2006). \textit{Watershed Models}. CRC Press, Boca Raton, FL.

\item Beven, K.J. (2012). \textit{Rainfall-Runoff Modelling: The Primer}, 2nd ed. Wiley-Blackwell,
Chichester, UK.
\end{enumerate}

\end{document}
